\documentclass[11pt, a4paper]{article}

% Gemeinsame Style-Definitionen (TEXINPUTS muss templates/ enthalten --- siehe scripts/build_tex.py)
% bfw-style.tex — gemeinsame Style-Definitionen für alle Handouts/Aufgabenhefte
% Voraussetzung: Kompilation mit xelatex (wegen fontspec)

\usepackage[ngerman]{babel}
\usepackage{geometry}
\geometry{a4paper, top=2.2cm, bottom=2.2cm, left=2.3cm, right=2.3cm}

\usepackage{fontspec}
% Fallback-Kette: Source Sans Pro -> Calibri -> Segoe UI -> Latin Modern Sans
% (Latin Modern ist immer mit MiKTeX vorhanden)
\IfFontExistsTF{Source Sans Pro}{%
  \setmainfont{Source Sans Pro}[Ligatures=TeX]%
}{\IfFontExistsTF{Calibri}{%
  \setmainfont{Calibri}[Ligatures=TeX]%
}{\IfFontExistsTF{Segoe UI}{%
  \setmainfont{Segoe UI}[Ligatures=TeX]%
}{%
  \setmainfont{Latin Modern Sans}[Ligatures=TeX]%
}}}
\IfFontExistsTF{Source Code Pro}{%
  \setmonofont{Source Code Pro}[Scale=0.92]%
}{\IfFontExistsTF{Consolas}{%
  \setmonofont{Consolas}[Scale=0.92]%
}{%
  \setmonofont{Latin Modern Mono}[Scale=0.92]%
}}

\usepackage{xcolor}
% Farbpalette: hell, freundlich, modern
\definecolor{bfwPrimary}{HTML}{0EA5A4}    % Petrol/Türkis - Primär
\definecolor{bfwPrimaryDark}{HTML}{0F766E} % dunkler Petrol für Headings
\definecolor{bfwAccent}{HTML}{F59E0B}     % Amber - Akzent
\definecolor{bfwSuccess}{HTML}{10B981}    % Grün - Erfolg / Lösung
\definecolor{bfwWarn}{HTML}{F97316}       % Orange - Achtung
\definecolor{bfwInk}{HTML}{0F172A}        % fast Schwarz
\definecolor{bfwInkSoft}{HTML}{475569}    % gedämpftes Grau
\definecolor{bfwBgSoft}{HTML}{F8FAFC}     % off-white
\definecolor{bfwBgTeal}{HTML}{ECFEFF}     % helles Türkis
\definecolor{bfwBgAmber}{HTML}{FEF3C7}    % helles Gelb
\definecolor{bfwBgRose}{HTML}{FFE4E6}     % helles Rosé für Eselsbrücken

\usepackage[colorlinks=true, linkcolor=bfwPrimaryDark, urlcolor=bfwPrimaryDark]{hyperref}
\usepackage{enumitem}
\setlist[itemize]{leftmargin=*, itemsep=2pt, topsep=2pt}
\setlist[enumerate]{leftmargin=*, itemsep=2pt, topsep=2pt}

\usepackage[most]{tcolorbox}
\usepackage{booktabs}
\usepackage{tikz}
\usetikzlibrary{decorations.pathmorphing, positioning, shapes.geometric, arrows.meta}

% Merkkästchen — wichtige Definition / Kernpunkt
\newtcolorbox{merksatz}[1][]{%
  enhanced, breakable,
  colback=bfwBgTeal,
  colframe=bfwPrimary,
  boxrule=0.6pt,
  arc=4pt,
  left=10pt, right=10pt, top=8pt, bottom=8pt,
  fonttitle=\bfseries\color{bfwPrimaryDark},
  title={\faLightbulb\ Merksatz},
  #1
}

% Eselsbrücke — Mnemoniks
\newtcolorbox{eselsbruecke}[1][]{%
  enhanced, breakable,
  colback=bfwBgRose,
  colframe=bfwAccent,
  boxrule=0.6pt,
  arc=4pt,
  left=10pt, right=10pt, top=8pt, bottom=8pt,
  fonttitle=\bfseries\color{bfwWarn},
  title={Eselsbrücke},
  #1
}

% Beispiel-Box
\newtcolorbox{beispiel}[1][]{%
  enhanced, breakable,
  colback=bfwBgSoft,
  colframe=bfwInkSoft,
  boxrule=0.4pt,
  arc=3pt,
  left=10pt, right=10pt, top=8pt, bottom=8pt,
  fonttitle=\bfseries\color{bfwInk},
  title={Beispiel},
  #1
}

% Lösungs-Box (für Aufgabenhefte)
\newtcolorbox{loesung}[1][]{%
  enhanced, breakable,
  colback=bfwBgAmber!40,
  colframe=bfwSuccess,
  boxrule=0.6pt,
  arc=3pt,
  left=10pt, right=10pt, top=8pt, bottom=8pt,
  fonttitle=\bfseries\color{bfwSuccess!70!black},
  title={Lösung},
  #1
}

% Glossar-Eintrag
\newcommand{\glossareintrag}[2]{%
  \par\noindent\textbf{\color{bfwPrimaryDark}#1} \enspace #2\par\smallskip
}

% Section-Stil — etwas farbiger als Standard
\usepackage{titlesec}
\titleformat{\section}
  {\Large\bfseries\color{bfwPrimaryDark}}
  {\thesection}{0.6em}{}
\titleformat{\subsection}
  {\large\bfseries\color{bfwPrimary}}
  {\thesubsection}{0.5em}{}
\titleformat{\subsubsection}
  {\normalsize\bfseries\color{bfwInk}}
  {\thesubsubsection}{0.4em}{}

% TikZ-Stil "skizzenhaft" — etwas weicher als Rough.js, aber stabil
\tikzset{
  roughbox/.style = {
    draw=bfwPrimaryDark, thick, fill=bfwBgTeal,
    rounded corners=4pt, inner sep=6pt
  },
  roughaccent/.style = {
    draw=bfwAccent, thick, fill=bfwBgAmber,
    rounded corners=4pt, inner sep=6pt
  },
  roughline/.style = {
    draw=bfwInkSoft, thick, -{Stealth[length=2.5mm]}
  }
}

% Bullet-Symbole (bewusst kein FontAwesome — vermeidet Font-Installations-Probleme)
\newcommand{\faLightbulb}{$\bullet$}
\newcommand{\faBookmark}{$\bullet$}

% Header/Footer
\usepackage{fancyhdr}
\pagestyle{fancy}
\fancyhf{}
\renewcommand{\headrulewidth}{0pt}
\fancyfoot[L]{\small\color{bfwInkSoft}\thepage}
\fancyfoot[R]{\small\color{bfwInkSoft} BFW M-064 Beschaffung im Büromanagement}


\title{\vspace*{-1.5cm}\Huge\bfseries\color{bfwPrimaryDark}Tag~2: Büroausstattung \& ergonomischer Bildschirmarbeitsplatz\\[0.4em]
  \large\color{bfwInkSoft}\textnormal{Die Arbeit an den Menschen anpassen}}
\author{\textsf{Büroprozesse gestalten und Arbeitsvorgänge organisieren \textperiodcentered{} M-067}\\
\textsf{\small\copyright\ Can Siebert\,\textperiodcentered\,2026}}
\date{30.06.2026}

\begin{document}
\maketitle
\thispagestyle{empty}

\vspace{1em}
\begin{merksatz}[title={\bfseries Worum geht es heute?}]
Gestern haben wir den Raum geplant --- heute richten wir den \emph{Arbeitsplatz} darin ein.
Auf dem Bürostuhl verbringen Büroangestellte den größten Teil ihrer Arbeitszeit. Falsche
Stühle, zu niedrige Bildschirme und spiegelnde Fenster verursachen Rücken-, Nacken- und
Kopfschmerzen --- Muskel- und Skeletterkrankungen sind die Hauptursache für Krankheitstage.
In diesem Handout lernen Sie, wie ein Bildschirmarbeitsplatz ergonomisch ausgestattet und
eingestellt wird, was die Arbeitsstättenverordnung verlangt und woran Sie geprüfte
Produkte erkennen. Das ist Grundlagenwissen für Gesundheitsschutz und IHK-Prüfung.
\end{merksatz}

\tableofcontents
\vspace{1em}
\hrule
\vspace{1em}

\section{Ergonomie --- die Arbeit an den Menschen anpassen}

Die \emph{Ergonomie} ist ein Teilgebiet der Arbeitswissenschaft. Sie befasst sich mit der
Abstimmung zwischen Mensch, Maschine und Arbeitswelt und dient dem \emph{präventiven
Gesundheitsschutz}. Dabei geht es vor allem darum, die Arbeit an den Menschen anzupassen ---
nicht den Menschen an die Arbeit. Dazu gehören die körpergerechte Gestaltung des
Arbeitsplatzes ebenso wie die Humanisierung der Arbeit. Ziel der Ergonomie ist es, die
arbeitenden Menschen vor Berufskrankheiten, Arbeitsunfällen und Monotonie zu bewahren.
Verwendete Arbeitsmittel wie Stühle, Tische oder Türgriffe sollen daher an die Maße des
menschlichen Körpers angepasst werden, um den Nutzer möglichst wenig zu belasten.

\begin{merksatz}
\textbf{Ergonomie} = die Arbeit an den Menschen anpassen. Ziel: Schutz vor
Berufskrankheiten, Arbeitsunfällen und Monotonie.
\end{merksatz}

\section{Der Bürostuhl}

Auf dem Bürostuhl verbringen Büroangestellte den größten Teil ihrer Arbeitszeit. Wird diese
Zeit auf ungeeigneten Stühlen oder in falscher Sitzposition verbracht, drohen Rücken-,
Kopfschmerzen und Verspannungen. Deshalb müssen sich Bürostühle an den jeweiligen Nutzer
anpassen lassen und bestimmte Mindestanforderungen erfüllen. Die europäische Norm
\emph{DIN EN 1335} regelt die grundlegenden Maße, Sicherheitsbestimmungen und Prüfkriterien
für Bürodrehstühle.

\subsection{Anforderungen nach DIN EN 1335}

\begin{itemize}
  \item leichte Federung beim Setzen, keine scharfen Kanten
  \item gepolsterte, atmungsaktive Sitzflächen- und Rückenlehnen
  \item stufenlos verstellbare Sitzhöhe von 42--50 cm
  \item Sitztiefe mindestens 38--44 cm, Sitzbreite mindestens 40--48 cm
  \item stufenlos verstellbare Rückenlehne mit einer Breite von mindestens 36--48 cm
  \item stand- und kippsicher, mindestens 5 wegrollsichere Rollen
  \item Rollwiderstand passend zum Bodenbelag (andere Rollen bei Teppich als bei glatten Böden)
  \item vorhandene Armlehnen in der Höhe verstellbar
\end{itemize}

Erfüllt ein Stuhl diese Anforderungen, muss er noch korrekt auf den Nutzer eingestellt
werden, damit dieser richtig sitzt.

\section{So sitzen Sie richtig --- der Bildschirmarbeitsplatz}

Erst der richtig eingestellte Stuhl, dann der Tisch, dann der Bildschirm: In dieser
Reihenfolge wird ein Bildschirmarbeitsplatz eingerichtet. Die folgenden fünf Regeln fassen
die ergonomisch korrekte Haltung zusammen.

\begin{enumerate}
  \item Die oberste Bildschirmzeile soll \emph{leicht unterhalb der waagerechten Sehachse}
    liegen --- der Blick ist also leicht gesenkt.
  \item Für den Monitor gilt ein \emph{Sichtabstand von mindestens 50 cm} (praktisch
    50--70 cm). Der Bildschirm soll im rechten Winkel zum Fenster stehen.
  \item Tastatur und Maus befinden sich in einer Ebene mit Ellenbogen und Handflächen.
  \item Ein \emph{90\textdegree-Winkel} zwischen Ober- und Unterarm sowie zwischen Ober-
    und Unterschenkel.
  \item Die Füße benötigen eine feste Auflage --- gegebenenfalls einen Fußhocker nutzen.
\end{enumerate}

\begin{eselsbruecke}
\textbf{„50 sehen, 90 beugen, fest stehen."} --- 50 cm Blickabstand, 90\textdegree-Winkel an
Armen und Beinen, feste Fußauflage. Drei Zahlen, der ganze Arbeitsplatz.
\end{eselsbruecke}

\subsection{Dynamisches Sitzen}

Um Haltungsschäden vorzubeugen, sollte man auch bei einem korrekt eingestellten Stuhl die
Sitzposition immer wieder verändern. Das nennt man \emph{dynamisches Sitzen}. Dabei werden
drei Sitzpositionen unterschieden:

\begin{itemize}
  \item \textbf{Vordere} Position: Winkel zwischen Oberkörper und Horizontale unter
    90\textdegree, die Oberschenkel sind nach unten geneigt.
  \item \textbf{Mittlere} Position: aufrechtes Sitzen, Oberkörper und Horizontale bilden
    einen 90\textdegree-Winkel.
  \item \textbf{Hintere} Position: zurückgelehnt, der Winkel ist größer als 90\textdegree.
\end{itemize}

\subsection{Steh-Sitz-Dynamik}

Wichtig ist, dass neben dem dynamischen Sitzen auch regelmäßig ein Wechsel zwischen Sitzen
und Stehen stattfindet. Die \emph{Steh-Sitz-Dynamik} verfolgt den Ansatz, dass zwei bis vier
Haltungswechsel pro Stunde erfolgen. Die Stehphasen sollten nicht länger als 20 bis 30
Minuten sein. Schweizer Arbeitsmediziner empfehlen die Faustregel:

\begin{center}
\begin{tabular}{l r}
\toprule
\textbf{Haltung} & \textbf{Anteil der Arbeitszeit}\\
\midrule
Sitzen & 60 \%\\
Stehen & 30 \%\\
Gezieltes Umhergehen & 10 \%\\
\bottomrule
\end{tabular}
\end{center}

\subsection{Sitzmechaniken bei Bürostühlen}

Die Neigbarkeit von Rückenlehne und Sitzfläche entlastet beim Bewegen die Bandscheiben.
Dafür gibt es vier verschiedene Sitzmechaniken:

\begin{center}
\begin{tabular}{p{3.2cm} p{9.2cm}}
\toprule
\textbf{Mechanik} & \textbf{Funktionsweise}\\
\midrule
Synchronmechanik & Neigung von Rückenlehne und Sitzfläche sind fest gekoppelt; beim Zurücklehnen ändert sich automatisch die Neigung der Sitzfläche und damit der Winkel zwischen Oberkörper und Oberschenkeln.\\
Wippmechanik & Beide Flächen sind gekoppelt, der Winkel zwischen ihnen bleibt aber konstant --- es entsteht ein Gefühl wie auf einem Schaukelstuhl.\\
Asynchronmechanik & Die Neigungswinkel von Sitzfläche und Rückenlehne lassen sich unabhängig voneinander verändern.\\
Pendelmechanik & Die Neigung des Stuhls wird durch Bewegung bzw.\ Gewichtsverlagerung gesteuert --- ähnlich wie auf einem Sitzball.\\
\bottomrule
\end{tabular}
\end{center}

\section{Prüfsiegel: GS und CE}

Um den Einkauf von Bürostühlen (und anderen Produkten) zu erleichtern, gibt es Prüfsiegel,
die die Einhaltung bestimmter Sicherheits- oder Qualitätskriterien anzeigen.

\begin{itemize}
  \item Das \textbf{GS-Zeichen} („Geprüfte Sicherheit") bescheinigt, dass ein Produkt die
    Anforderungen des Produktsicherheitsgesetzes erfüllt. DIN-Normen, Europäische Normen
    oder andere anerkannte Regeln der Technik konkretisieren diese Anforderungen. Die
    Unternehmen verwenden dieses Siegel \emph{freiwillig}.
  \item Mit dem \textbf{CE-Zeichen} (Communautés Européennes) bestätigt der Hersteller in
    Eigenverantwortung, dass sein Produkt die Anforderungen der EU erfüllt. Es ist ein
    \emph{gesetzlich vorgeschriebenes} Verwaltungszeichen --- etwas wie ein „technischer
    Reisepass", damit das Produkt frei im Binnenmarkt verkehren darf.
\end{itemize}

\begin{merksatz}
GS und CE sind die einzig gesetzlich geregelten Zeichen für Produktsicherheit in Europa.
\textbf{GS = freiwillig}, \textbf{CE = Pflicht}. Sonstige Siegel (z.\,B.\ TÜV) vergeben
private Prüfstellen.
\end{merksatz}

\section{Der Arbeitstisch}

Auch der Arbeitstisch muss an Körpergröße und Proportionen des Nutzers anpassbar sein. Man
unterscheidet \emph{einstellbare} und \emph{verstellbare} Tische: Ein einstellbarer Tisch
wird nur einmal mit Werkzeug auf einen Nutzer eingestellt; ein verstellbarer Tisch lässt
sich jederzeit mit wenigen Handgriffen in der Höhe anpassen. Die richtige Tischhöhe ist
erreicht, wenn bei geradem Rücken und korrekt eingestelltem Stuhl die Unterarme flach auf
der Platte aufliegen und die Schultern nicht hochgezogen werden müssen.

\subsection{Anforderungen an den Arbeitstisch}

\begin{itemize}
  \item Arbeitsfläche von mindestens 160 cm $\times$ 80 cm (Ausnahme z.\,B.\ Call-Center
    mit reiner PC-Arbeit: 120 cm $\times$ 80 cm)
  \item Mindesttiefe von 80 cm, genügend Beinfreiheit
  \item höheneinstellbar zwischen 65 cm und 85 cm (ansonsten mindestens 72 cm Höhe)
  \item abgerundete Kanten, reflexionsarme Oberfläche
  \item geringe Wärmeleitfähigkeit (Holz- oder Kunststoffplatten statt Glas oder Metall)
  \item ausreichend Stauraum für Unterlagen, flexible Anordnung von Maus und Tastatur
\end{itemize}

Neben dem reinen Sitzarbeitstisch gibt es \textbf{Steh-Sitz-Arbeitstische}, deren
Arbeitshöhe per Elektromotor oder Gasfeder zwischen 62 cm und 120 cm verändert werden kann,
sowie reine \textbf{Steharbeitsplätze} (Stehhöhe 95--120 cm). Ein \emph{Drei-Zonen-Arbeitsplatz}
vereint alle Tischarten: eine Zone für reine Sitzarbeit ohne PC, eine für Sitzarbeit am PC
und eine für Arbeit im Stehen.

\begin{beispiel}
Frau Akin richtet ihren neuen Arbeitsplatz ein. Zuerst stellt sie den Bürostuhl so ein,
dass Ober- und Unterschenkel einen 90\textdegree-Winkel bilden und die Füße fest stehen.
Dann passt sie den höhenverstellbaren Tisch so an, dass ihre Unterarme flach aufliegen.
Den externen Bildschirm rückt sie auf rund 60 cm Abstand und dreht ihn im rechten Winkel
zum Fenster --- die Spiegelungen verschwinden. Tastatur und Maus legt sie auf
Ellenbogenhöhe. Zweimal pro Stunde fährt sie den Tisch hoch und arbeitet im Stehen.
\end{beispiel}

\section{Komponenten und die Arbeitsstättenverordnung}

Der \textbf{Bildschirmarbeitsplatz} besteht neben Bürostuhl und Arbeitstisch in der Regel
aus folgenden Komponenten: \emph{Bildschirm, Tastatur, Maus, Vorlagenhalter, Software und
Drucker}.

Die \emph{Arbeitsstättenverordnung (ArbStättV)} legt fest, was der Arbeitgeber beim
Einrichten und Betreiben von Arbeitsstätten in Bezug auf Sicherheit und Gesundheitsschutz
beachten muss --- z.\,B.\ Anforderungen an Arbeitsräume, Beleuchtung, Belüftung oder
Raumtemperatur. Sie enthält Schutzziele und allgemein gehaltene Anforderungen; konkretisiert
werden diese in den \emph{Technischen Regeln für Arbeitsstätten (ASR)}. Im Anhang
„Anforderungen und Maßnahmen für Arbeitsstätten nach § 3 Absatz 1" werden unter Punkt 6 die
Anforderungen an den Bildschirmarbeitsplatz konkretisiert (ehemals Bildschirmarbeitsverordnung).

\subsection{Wichtige Anforderungen an Bildschirmarbeitsplätze}

\begin{itemize}
  \item Bildschirmarbeitsplätze so einrichten, dass Sicherheit und Gesundheit gewährleistet
    sind; Grundsätze der Ergonomie anwenden.
  \item Für \emph{ausreichend Raum für wechselnde Arbeitshaltungen und -bewegungen} sorgen;
    Tätigkeiten durch andere Tätigkeiten oder Erholungszeiten unterbrechen.
  \item Bildschirme so aufstellen, dass die Oberflächen frei von störenden Reflexionen und
    Blendungen sind; reflexionsarme, dreh- und neigbare Oberflächen.
  \item Arbeitstische/-flächen mit reflexionsarmer Oberfläche; vor der Tastatur muss ein
    Auflegen der Handballen möglich sein.
  \item Auf Wunsch \emph{Fußstütze und Manuskripthalter} bereitstellen.
  \item \emph{Beleuchtung} der Arbeitsaufgabe und dem Sehvermögen anpassen; angemessener
    Kontrast, keine Blendungen oder Spiegelungen.
  \item Bild \emph{flimmerfrei und verzerrungsfrei}; Helligkeit und Kontrast individuell
    einstellbar; Zeichengröße und Zeilenabstand anpassbar.
  \item \emph{Tastaturen} müssen vom Bildschirm getrennte, neigbare Einheiten mit
    reflexionsarmer Oberfläche sein; die Beschriftung muss gut lesbar sein.
  \item \emph{Tragbare Bildschirmgeräte} (Laptops) ohne Trennung von Bildschirm und Tastatur
    dürfen nur kurzzeitig betrieben werden --- im Dauereinsatz: Dockingstation mit externem
    Bildschirm, Tastatur und Maus.
  \item \emph{Softwareergonomie}: Die Software muss an Kenntnisse und Arbeitsaufgabe
    anpassbar sein, Angaben zu Dialogabläufen machen und Fehlerbeseitigung erlauben; eine
    heimliche Leistungs-/Verhaltenskontrolle ist unzulässig.
\end{itemize}

\section{Modernes Wissen für die Praxis}

\begin{itemize}
  \item \textbf{Steh-Sitz-Dynamik im Alltag}: Telefonate im Stehen führen, den Drucker
    bewusst weiter weg stellen --- so entstehen automatisch Haltungswechsel.
  \item \textbf{Laptop + Dockingstation}: Der Laptop allein ist ergonomisch ein Kompromiss
    (Bildschirm hoch \emph{oder} Tastatur tief). Mit Docking, externem Monitor, Tastatur und
    Maus lässt sich jedes Element korrekt positionieren.
  \item \textbf{Homeoffice-Ergonomie}: Auch zu Hause gelten dieselben Regeln. Notlösungen
    (Küchentisch, Sofa) führen schnell zu Beschwerden; ein höhenverstellbarer Stuhl, ein
    externer Bildschirm und gutes Licht sind die wichtigsten Investitionen.
\end{itemize}

\section{Zusammenfassung}

\begin{enumerate}
  \item \textbf{Ergonomie} passt die Arbeit an den Menschen an und schützt vor
    Berufskrankheiten, Unfällen und Monotonie.
  \item Der \textbf{Bürostuhl} muss DIN EN 1335 erfüllen (verstellbare Sitzhöhe 42--50 cm,
    mind. 5 Rollen, höhenverstellbare Armlehnen u.\,a.) und korrekt eingestellt werden.
  \item „So sitzen Sie richtig": Bildschirm leicht unter der Sehachse, Sichtabstand
    $\geq$ 50 cm, \textbf{90\textdegree-Winkel}, feste Fußauflage.
  \item \textbf{Dynamisches Sitzen} und \textbf{Steh-Sitz-Dynamik} (60/30/10, Wechsel
    2--4$\times$ pro Stunde); vier Sitzmechaniken.
  \item \textbf{GS} (freiwillig) und \textbf{CE} (Pflicht) sind die gesetzlich geregelten
    Produktsicherheitszeichen.
  \item Der \textbf{Arbeitstisch} (mind.\ 160 $\times$ 80 cm, höhenverstellbar 65--85 cm)
    und die \textbf{Arbeitsstättenverordnung} regeln den Bildschirmarbeitsplatz.
\end{enumerate}

\newpage

\section*{Glossar}
\addcontentsline{toc}{section}{Glossar}

\glossareintrag{Ergonomie}{Teilgebiet der Arbeitswissenschaft; passt die Arbeit an den Menschen an, um ihn vor Berufskrankheiten, Unfällen und Monotonie zu schützen.}

\glossareintrag{DIN EN 1335}{Europäische Norm mit den grundlegenden Maßen, Sicherheitsbestimmungen und Prüfkriterien für Bürodrehstühle.}

\glossareintrag{Bürodrehstuhl}{Bürostuhl, der sich an den Nutzer anpassen lässt (Sitzhöhe, Lehne, Armlehnen) und mind.\ 5 wegrollsichere Rollen hat.}

\glossareintrag{Bildschirmarbeitsplatz}{Arbeitsplatz aus Bürostuhl, Arbeitstisch, Bildschirm, Tastatur, Maus, Vorlagenhalter, Software und Drucker.}

\glossareintrag{Sichtabstand}{Empfohlener Abstand Auge--Bildschirm von mindestens 50 cm (praktisch 50--70 cm).}

\glossareintrag{Dynamisches Sitzen}{Häufiger Wechsel zwischen vorderer, mittlerer und hinterer Sitzposition zur Entlastung der Bandscheiben.}

\glossareintrag{Steh-Sitz-Dynamik}{Regelmäßiger Wechsel zwischen Sitzen und Stehen (2--4 Wechsel/Stunde); Faustregel 60 \% sitzen, 30 \% stehen, 10 \% gehen.}

\glossareintrag{Synchronmechanik}{Sitzmechanik, bei der Neigung von Sitzfläche und Rückenlehne fest gekoppelt sind.}

\glossareintrag{Asynchronmechanik}{Sitzmechanik, bei der die Neigungswinkel von Sitzfläche und Rückenlehne unabhängig verstellbar sind.}

\glossareintrag{GS-Zeichen}{„Geprüfte Sicherheit"; freiwilliges Siegel für die Einhaltung des Produktsicherheitsgesetzes.}

\glossareintrag{CE-Zeichen}{Gesetzlich vorgeschriebenes Zeichen; der Hersteller bestätigt die Einhaltung der EU-Anforderungen.}

\glossareintrag{Steh-Sitz-Arbeitstisch}{Tisch, dessen Arbeitshöhe per Elektromotor oder Gasfeder zwischen 62 cm und 120 cm verstellbar ist.}

\glossareintrag{Drei-Zonen-Arbeitsplatz}{Arbeitsplatz mit drei Zonen: Sitzarbeit ohne PC, Sitzarbeit am PC, Arbeit im Stehen.}

\glossareintrag{Arbeitsstättenverordnung (ArbStättV)}{Verordnung zu Sicherheit und Gesundheitsschutz beim Einrichten und Betreiben von Arbeitsstätten; Anhang Punkt 6 regelt Bildschirmarbeitsplätze.}

\glossareintrag{ASR}{Technische Regeln für Arbeitsstätten, die die Anforderungen der ArbStättV konkretisieren.}

\glossareintrag{Softwareergonomie}{Anpassbarkeit der Software an Kenntnisse und Arbeitsaufgabe; Unterstützung bei Dialogabläufen und Fehlerbeseitigung.}

\end{document}
